\section{Linear-strand growth for the two truncations}
\label{sec:growth}

Section~\ref{sec:betti-euler-shifts} studies $n\to\infty$ with monomial
degree $d$ fixed. Here we fix the homological degree $i$ and relation
arity $k$, sum over all Goldbach degrees $n$, and let the cutoff
$m\to\infty$.
For the binary ring the first possible syzygies in this homological degree
have monomial degree $i+1$. In the Koszul complex their coefficients
have monomial degree one. Their total number has a different leading constant for
each of our two truncations. Both constants can be computed by comparing
three consecutive Koszul terms over the given field $\kk$.

\subsection{An elementary estimate in the binary case}

Let $G=R/I$ be the binary Goldbach ring on a finite nonempty support
$\A$. In homological degree $i\ge1$, the linear strand contributes
$\beta_{i,*,i+1}$ to the total $\beta_{i,*,*}$.

\begin{proposition}\label{grow:binary-bound}
For every $i\ge1$,
\[
 \begin{aligned}
 \beta_{i,*,i+1}
 &\le |\A|\binom{|\A|}{i}-\binom{|\A|}{i+1},\\
 \beta_{i,*,i+1}
 &\ge |\A|\binom{|\A|}{i}-\binom{|\A|}{i+1}
       -|\A+\A|\binom{|\A|}{i-1}.
 \end{aligned}
\]
Consequently, along any family with $|\A|\to\infty$ and $|\A+\A|=o(|\A|^2)$,
\[
 \beta_{i,*,i+1}\sim\frac{i}{(i+1)!}|\A|^{i+1}
 \qquad\text{for every fixed }i\ge1.
\]
\end{proposition}
\begin{proof}
Write $G_{*,r}=\bigoplus_n G_{n,r}$, where $*$ denotes the direct sum
over all Goldbach degrees. In monomial degree $i+1$, the relevant
Koszul terms are
\[
 G_{*,0}\otimes\bigwedge^{i+1}\kk^{\A}
 \xrightarrow{\mathrm{d}_{i+1}}
 G_{*,1}\otimes\bigwedge^i\kk^{\A}
 \xrightarrow{\mathrm{d}_i}
 G_{*,2}\otimes\bigwedge^{i-1}\kk^{\A}.
\]
Their dimensions are $\binom{|\A|}{i+1}$, $|\A|\binom{|\A|}{i}$, and
$|\A+\A|\binom{|\A|}{i-1}$. Indeed, there are no linear relations, and
all quadratic monomials of a given represented weight agree in $G$.
The columns of $\mathrm{d}_{i+1}$ have disjoint supports and unit entries:
a term $x_a\varepsilon_{\sigma}$, with $a\notin \sigma$, determines its unique
source $\varepsilon_{\sigma\cup\{a\}}$. This map is therefore injective
over every field. Summing \eqref{eq:koszul-rank-nullity} over Goldbach
degrees $n$ therefore gives
\[
 \beta_{i,*,i+1}
 =|\A|\binom{|\A|}{i}-\binom{|\A|}{i+1}
   -\dim_\kk\operatorname{im}\mathrm{d}_i.
\]
Bounding the rank by zero and by the target dimension proves the
inequalities. For fixed $i$, the main term is asymptotic to
$i|\A|^{i+1}/(i+1)!$, while the error is $O_i(|\A+\A|\,|\A|^{i-1})$.
\end{proof}

Let $m$ denote the prime cutoff as in Section~\ref{setup:resolution-construction}.
For $\A_m=\{p\le m:p\text{ odd prime}\}$, the prime number theorem gives
$|\A_m|\sim m/\log m$.
Since $|\A_m+\A_m|=O(m)=o(|\A_m|^2)$, the proposition gives the binary
variable-cutoff asymptotic. More generally, the same proof applies to
any standard graded quotient without linear relations, with $|\A+\A|$
replaced by its degree-two Hilbert function.

\subsection{The linear strand at the two cutoffs}

Already among the cubic second syzygies the distinction is visible.
Summed over all weights, their number is asymptotic to $|\A_m|^3/3$
for $G^{<m}$ and to $|\A_m|^3/8$ for $G^{\le m}$. The same comparison
gives the constants in every fixed homological degree.

\begin{theorem}\label{grow:both-cutoffs}
For the binary ring over any field and every fixed $i\ge1$, as $m\to\infty$,
\[
 \beta^{<m}_{i,*,i+1}\sim\frac{i}{(i+1)!}|\A_m|^{i+1},
 \qquad
 \beta^{\le m}_{i,*,i+1}
       \sim\frac{1-2^{-i}}{(i+1)!}|\A_m|^{i+1}.
\]
More generally, fix $k\ge2$ and $i\ge1$, and let $G$ be the $k$-ary Goldbach
ring on the odd primes. Its variable cutoff $G^{<m}$ uses all
defining relations on $\A_m$, and its relation cutoff $G^{\le m}$
uses those of common weight at most $m$. Then
\begin{align}
 \beta^{<m}_{i,*,i+k-1}
 &\sim\frac{|\A_m|^{k+i-1}}{(k-1)!(i-1)!(k+i-1)},
 \label{grow:variable-constant}\\
 \beta^{\le m}_{i,*,i+k-1}
 &\sim C_{k,i}|\A_m|^{k+i-1},\qquad
 C_{k,i}=\frac{\displaystyle\sum_{j=1}^k
       (-1)^{j+1}\binom kj j^{1-i}}{k!(k+i-1)!}.
 \label{grow:relation-constant}
\end{align}
\end{theorem}

The next lemma uses only the monomial grading. Within its statement
$\beta_{i,d}(M)$ denotes $\dim_\kk\Tor_i^R(M,\kk)_d$.
For a bigraded quotient this is $\beta_{i,*,d}(M)$.
In its proof, $K_j(M)_d$ denotes the monomial-degree-$d$ part of the
$j$th Koszul term.

\begin{lemma}\label{grow:koszul-comparison}
Suppose $\mathfrak a\subseteq\mathfrak b\subseteq R=\kk[x_1,\ldots,x_N]$ are homogeneous
ideals with no elements in degrees below $k\ge2$. If
$s=\dim_\kk(\mathfrak b/\mathfrak a)_k$, then for $i\ge1$,
\[
 0\le\beta_{i,i+k-1}(R/\mathfrak b)-\beta_{i,i+k-1}(R/\mathfrak a)
       \le s\binom N{i-1}.
\]
\end{lemma}
\begin{proof}
Put $E=\kk^N$, $d=i+k-1$, and $K(M)=\mathbb K_{\kk}\otimes_R M$.
Since $K_j(M)_d=M_{d-j}\otimes_\kk\bigwedge^jE$, the coefficient
degrees in homological degrees $i+1,i,i-1$ are $k-2,k-1,k$.
The two quotients agree below degree $k$, so they have the same
incoming map $A:U\to V$, where
\[
 U=R_{k-2}\otimes_\kk\bigwedge^{i+1}E,\qquad
 V=R_{k-1}\otimes_\kk\bigwedge^iE.
\]
The common term $\dim_\kk(V/\operatorname{im}A)$ in
\eqref{eq:koszul-rank-nullity} cancels when the two Betti numbers are
subtracted. Only the outgoing targets differ:
\[
 W_{\mathfrak a}=(R/\mathfrak a)_k\otimes_\kk\bigwedge^{i-1}E,
 \qquad
 W_{\mathfrak b}=(R/\mathfrak b)_k\otimes_\kk\bigwedge^{i-1}E.
\]
The quotient $\rho:W_{\mathfrak a}\twoheadrightarrow W_{\mathfrak b}$
satisfies $B_{\mathfrak b}=\rho B_{\mathfrak a}$ for the outgoing maps,
and $\dim_\kk\ker\rho=s\binom N{i-1}$. Thus
$\ker B_{\mathfrak a}\subseteq\ker B_{\mathfrak b}$, while
$\operatorname{im}A$ is unchanged. The map $v\mapsto B_{\mathfrak a}v$
identifies
\[
 \ker B_{\mathfrak b}/\ker B_{\mathfrak a}
 \cong\operatorname{im}B_{\mathfrak a}\cap\ker\rho.
\]
Consequently
\[
 0\le\beta_{i,d}(R/\mathfrak b)-\beta_{i,d}(R/\mathfrak a)
 =\dim_\kk(\operatorname{im}B_{\mathfrak a}\cap\ker\rho)
 \le s\binom N{i-1}.
\]
\end{proof}

\begin{proof}[Proof of Theorem~\ref{grow:both-cutoffs}]
Within this proof, write $I$ for the binomial ideal of the chosen
truncation and $I'$ for its monomial comparison ideal. For the relation
cutoff, take $I=I^{\le m}$ and
\[
 I'=\bigl(x_{p_1}\cdots x_{p_k}\mid p_1,\ldots,p_k\in\A_m,
       \ \textstyle\sum_{j=1}^k p_j\le m\bigr).
\]
Both ideals have no elements below monomial degree $k$, and
$I\subseteq I'$ because each generator of $I$ is a difference of two
listed monomials.

In monomial degree $k$, the binomial ideal $I$ is the coefficient-sum-zero
subspace of each represented fiber of weight at most $m$, whereas $I'$
contains the whole fiber. Therefore $s=\dim_\kk(I'/I)_{*,k}$ counts the
represented weights and is $O_k(m)$. For the variable cutoff, take instead
$I=I^{<m}$ and $I'=\m^k$, where $\m=(x_p\mid p\in\A_m)$;
then $s=|k\A_m|=O_k(m)$. In either case the error from
Lemma~\ref{grow:koszul-comparison} is
\[
 O_{k,i}(m |\A_m|^{i-1})=o(|\A_m|^{k+i-1}).
\]
Indeed, after division by $|\A_m|^{k+i-1}$, the error is
$O_{k,i}(m/|\A_m|^k)$, and the prime number theorem gives
\[
 \frac{m}{|\A_m|^k}
 \sim\frac{(\log m)^k}{m^{k-1}}\longrightarrow0.
\]

List the variables in increasing order of prime label. Both comparison
ideals are strongly stable: replacing a variable by one with smaller
index does not increase the additive weight. They are generated in
degree $k$. If $q(x^\alpha)$ is the largest index of a variable dividing
a generator $x^\alpha$, the Eliahou--Kervaire formula \cite{EK90} gives
\begin{equation}\label{grow:EK}
 \beta_{i,*,i+k-1}(R_m/I')
 =\sum_{x^\alpha\in\mathcal G(I')}\binom{q(x^\alpha)-1}{i-1}.
\end{equation}
All positive Betti numbers of the comparison quotient $R_m/I'$ lie on
this strand. The formula is valid
over every field. In the case $I'=\m^k$, the sum becomes
\[
 \sum_{j=1}^{|\A_m|}\binom{j+k-2}{k-1}\binom{j-1}{i-1}
 \sim\frac{|\A_m|^{k+i-1}}{(k-1)!(i-1)!(k+i-1)}.
\]

For the relation cutoff, the prime number theorem gives
$|\A_{tm}|/|\A_m|\to t$ for each fixed $t\in(0,1]$.
Monotonicity and continuity of the limiting distribution function make
this convergence uniform on $[0,1]$.
Thus the measures $|\A_m|^{-1}\sum_{p\in\A_m}\delta_{p/m}$ converge weakly
to uniform measure, and $q(x^\alpha)/|\A_m|$ differs uniformly by $o(1)$ from
the largest scaled label in $x^\alpha$. Generators with a repeated label
contribute only $O_{k,i}(|\A_m|^{k+i-2})$ to \eqref{grow:EK}. Counting the
remaining ordered tuples and dividing by $k!$ therefore gives
\begin{equation}\label{grow:simplex-moment}
 \frac{\sum_n\beta_{i,n,i+k-1}(R_m/I')}{|\A_m|^{k+i-1}}
 \longrightarrow\frac{1}{k!(i-1)!}
 \int_{\substack{u_j\ge0\\\sum_ju_j\le1}}
       (\max_j u_j)^{i-1}\,du.
\end{equation}
The simplex boundary has measure zero, so the cutoff indicator presents
no obstruction to this limit.

Put $r=i-1$. For $r\ge1$, inclusion--exclusion yields
\[
 \operatorname{vol}\{u\ge0:\textstyle\sum_ju_j\le1,\ \max_j u_j>s\}
 =\frac{1}{k!}\sum_{j=1}^k(-1)^{j+1}\binom kj(1-js)_+^k.
\]
Integrating against $r s^{r-1}\,ds$ and evaluating the beta integral gives
\[
 \int_{\sum_ju_j\le1}(\max_j u_j)^r\,du
 =\frac{r!}{(k+r)!}\sum_{j=1}^k(-1)^{j+1}\binom kj j^{-r}.
\]
For $r=0$ the same formula holds, because the simplex has volume $1/k!$
and the alternating binomial sum is one. Substitution into
\eqref{grow:simplex-moment}, followed by the rank comparison, proves
\eqref{grow:relation-constant}. At $k=2$, the sum simplifies to the
stated binary constant.
\end{proof}

The ratio of the binary
relation-cutoff linear strand to the variable-cutoff linear strand
tends to $(1-2^{-i})/i$. These leading constants are independent of
the characteristic, although individual Betti numbers need not be.

\subsection{The first nonlinear strand}

The first nonlinear contribution to the binary variable cutoff is
smaller. For $\rho\ge0$, write uniquely
$\rho=\binom a2+b$ with $a\ge1$ and $0\le b<a$, and set
\[
 \rho^{\langle2\rangle}=\binom{a+1}{3}+\binom{b+1}{2}.
\]
This is the degree-two Macaulay successor, and
$\rho^{\langle2\rangle}=O(\rho^{3/2})$ for $\rho\ge1$.

\begin{proposition}\label{grow:nonlinear-bound}
Let $G=R/I$ be the binary ring on a finite nonempty support $\A$, and
put $\rho=|\A+\A|$.
Over every field and for every $i\ge1$,
\[
 \beta_{i,*,i+2}
 \le\binom{|\A|-1}{i-1}
       \bigl(\rho^{\langle2\rangle}-\dim_\kk G_{*,3}\bigr)
 \le\binom{|\A|-1}{i-1}\rho^{\langle2\rangle}.
\]
\end{proposition}
\begin{proof}
Write $H(d)=\dim_\kk G_{*,d}$. Order the support as
$a_1<\cdots<a_{|\A|}$, and let $L\subseteq R$ be the lex-segment ideal
whose quotient $R/L$ has Hilbert function $H$, using
$x_{a_1}>\cdots>x_{a_{|\A|}}$.
Macaulay's theorem supplies this ideal; see the discussion preceding
\cite[Lemmas~2--3]{JKV21}. It has no linear generators since $H(1)=|\A|$.

Let $Q$ be the ideal generated by the quadratic monomials of $L$.
Since $H(2)=\rho$ and the quadratic lex segment attains Macaulay's
growth bound,
\[
 \dim_\kk(R/Q)_{*,3}=\rho^{\langle2\rangle}.
\]
The cubic monomials of $L$ not divisible by its quadratic generators
form a basis of $(L/Q)_{*,3}$ and are precisely its minimal cubic
generators, denoted $\mathcal G(L)_3$. Hence
\[
 |\mathcal G(L)_3|=\dim_\kk(L/Q)_{*,3}
 =\rho^{\langle2\rangle}-H(3).
\]
These are generators of the comparison ideal $L$, even though $I$
itself is generated by quadrics.

For a monomial $x^\alpha$, put
$q(x^\alpha)=\max\{j:x_{a_j}\text{ divides }x^\alpha\}$.
The Bigatti--Hulett--Pardue theorem and the Eliahou--Kervaire formula,
valid over arbitrary fields, give
\[
 \beta_{i,*,i+2}\le\beta_{i,*,i+2}(R/L)
 =\sum_{x^\alpha\in\mathcal G(L)_3}\binom{q(x^\alpha)-1}{i-1}
 \le\binom{|\A|-1}{i-1}\bigl(\rho^{\langle2\rangle}-H(3)\bigr);
\]
see \cite[Lemmas~2--3]{JKV21} and \cite{EK90}.
The remaining bound follows from $H(3)\ge0$. Finally,
$a=O(\sqrt\rho)$ and $b<a$ in the Macaulay representation give
$\rho^{\langle2\rangle}=O(\rho^{3/2})$.
\end{proof}

For the odd primes at most $m$, the support has size $|\A_m|$ and
$\rho=|\A_m+\A_m|=O(m)$. As $m^{3/2}/|\A_m|^2\sim(\log m)^2/\sqrt m\to0$,
the proposition gives
\[
 \beta^{<m}_{i,*,i+2}
 =O_i(|\A_m|^{i-1}m^{3/2})=o(|\A_m|^{i+1}).
\]
Thus for the variable cutoff the sum of the first two strands has
the same leading asymptotic as the linear strand alone. Determining
the asymptotics of the total Betti number $\beta^{<m}_{i,*,*}$ requires
control of monomial degrees $i+3,i+4,\ldots$. All estimates here are for fixed $i$.
